\section{Protected fibers and off-fiber response}
\label{sec:protected-response}

Let $H(\boldsymbol\lambda)$ be a differentiable family of self-adjoint
Hamiltonians.  Throughout this paper, $E_0(\boldsymbol\lambda)$ denotes an
isolated eigenvalue of constant multiplicity $D$ on the parameter region under
consideration, and $P(\boldsymbol\lambda)$ is its constant rank $D$ spectral
projector.  We write
$Q=1-P$ and assume a nonzero spectral gap
$\Delta=\operatorname{dist}[E_0,\operatorname{spec}(QHQ)]$.  These assumptions
are stronger than observing a numerically narrow band: they ensure that the
projector is differentiable and that the inverse of $H-E_0$ on $Q\cH$ is
bounded.  No response formula below is continued through a point at which the
rank changes or the spectral gap closes.

\subsection{The splitting and mixing blocks}

For a coordinate $\lambda^a$, set $h_a=\partial_aH$.  The two blocks of $h_a$
that are relevant to the protected fiber are
\begin{equation}
 S_a=Ph_aP-\frac{\Tr(Ph_aP)}{D}P,
 \qquad
 V_a=Qh_aP .
 \label{eq:splitting}
\end{equation}
The Hermitian operator $S_a$ is the traceless intrafiber splitting channel,
whereas $V_a$ maps the protected fiber into its orthogonal complement.  If the
eigenvalue remains exactly $D$-fold degenerate along the path, differentiating
$HP=E_0P$ and projecting with $P$ on both sides gives
$Ph_aP=(\partial_aE_0)P$, hence $S_a=0$.  This is a local necessary condition
for an exactness-preserving path.  It is not, by itself, a claim that an
arbitrary linear perturbation with $S_a=0$ preserves the degeneracy at finite
amplitude; the finite path must continue to satisfy the rank and gap
conditions.

Crucially, $S_a=0$ imposes no corresponding condition on $V_a$.  Exact
degeneracy fixes the restriction of the Hamiltonian to $P\cH$, but it does not
fix how that subspace is embedded in the full Hilbert space.  This distinction
also works in the opposite direction.  An intervention supported entirely in
$P\cH$ can make $S_a$ nonzero, and can even produce random-matrix level
statistics inside the multiplet, while keeping the projector fixed.  Such an
intervention changes intrafiber energies but produces no quantum geometry.

\subsection{Reduced response of an isolated fiber}

Under the assumptions above, the reduced resolvent
\[
 R_Q=Q\bigl[Q(H-E_0)Q\bigr]^{-1}Q
\]
exists and satisfies $\lVert R_Q\rVert\leq\Delta^{-1}$.  Projecting the
derivative of $HP=E_0P$ from $P\cH$ to $Q\cH$ gives the off-fiber response
\begin{equation}
 X_a=Q(\partial_aP)P=-R_QV_a .
 \label{eq:response}
\end{equation}
Thus a large response can arise either from a large coupling matrix element
$V_a$ or from a small excitation denominator.  The distinction matters when
systems with different gaps are compared: an unnormalized response does not,
by itself, measure more effective mixing channels.  Equation~\eqref{eq:response}
also makes the domain of the calculation explicit.  It is an adiabatic
response of an isolated spectral subspace, not a prescription for following a
collection of levels after that subspace touches the rest of the spectrum.

The matrix-valued quantum geometric tensor is the operator on $P\cH$
\begin{equation}
 \mathcal Q_{ab}
 =P(\partial_aP)Q(\partial_bP)P
 =X_a^\dagger X_b,
 \qquad
 g_{ab}=\frac{\mathcal Q_{ab}+\mathcal Q_{ba}}{2}.
 \label{eq:qgt}
\end{equation}
The diagonal $g_{aa}=X_a^\dagger X_a$ is positive semidefinite, and its rank
counts the fiber directions resolved by tangent $a$, subject to the gap
weighting in $R_Q$.  Unlike the scalar quantum metric of a nondegenerate
state, the off-diagonal $g_{ab}$ is generally a matrix inside the protected
multiplet; no analogous rank interpretation is assigned to it.

Choose a local orthonormal frame $\Psi=(|u_1\rangle,\ldots,|u_D\rangle)$ for
$P\cH$.  With the canonical convention
$A_a=i\Psi^\dagger\partial_a\Psi$, the non-Abelian Berry curvature is
\begin{equation}
 \begin{split}
 F_{ab}&=\partial_aA_b-\partial_bA_a-i[A_a,A_b]\\
 &=i(\mathcal Q_{ab}-\mathcal Q_{ba})
 =i\,P[\partial_aP,\partial_bP]P .
 \end{split}
 \label{eq:curvature}
\end{equation}
For $D=1$, this convention reduces to
$F_{ab}=-2\operatorname{Im}\mathcal Q_{ab}$.  Equations~\eqref{eq:qgt} and
\eqref{eq:curvature} are therefore the symmetric and antisymmetric parts of
the same off-fiber response.

\subsection{Gauge covariance and observable contractions}

A change of frame $\Psi\mapsto\Psi U(\boldsymbol\lambda)$ leaves the projector
unchanged but sends
\begin{align}
 A_a&\mapsto U^\dagger A_aU+iU^\dagger\partial_aU,\nonumber\\
 F_{ab}&\mapsto U^\dagger F_{ab}U,
 &\mathcal Q_{ab}&\mapsto U^\dagger\mathcal Q_{ab}U.
\end{align}
The connection transforms inhomogeneously, while the curvature and QGT are
gauge covariant.  Raw entries of $A_a$ are consequently not primary
observables in our statistical tests.  The eigenvalues of a fixed Hermitian
curvature component, traces such as $\Tr F_{ab}^k$, and closed words such as
$\Tr(\mathcal Q_{ab}\mathcal Q_{cd})$ are gauge invariant.  When rectangular
response matrices are used computationally, their fiber-frame representation
transforms as $X_a\mapsto X_aU$; every reported contraction closes the fiber
indices and has the corresponding invariant or covariant transformation law.

This establishes the separation used throughout the paper.  The condition
$S_a=0$ protects the eigenvalue to first order, whereas $V_a$ and the reduced
resolvent determine the motion of the protected fiber.  A fixed projector has
$X_a=g_{ab}=F_{ab}=0$.  A moving projector can instead have scalar,
block-reducible, or fully resolved curvature.  Which of these cases displays
random-matrix correlations is a statistical question, not a consequence of
the degeneracy alone.
